% p.3 — فهرست المضامين (contents of the lithograph). Two columns on the page: right (1–20), left (21–35).
\mspage{3}
\gdef\currentbab{فهرست المضامين}
% Layout: as printed, two tables side by side (1–20 right, 21–35 left)
\begin{tables}
\begin{sarf}{|c|>{\raggedleft\arraybackslash}p{.33\linewidth}|c|}
\thead{نمبر} & \thead{مضامين} & \thead{صفحه}\\\hline
1 & الصرف الصغير & 4\\
2 & الصرف الكبير & 5\\
3 & المنفيات & 8\\
4 & المؤكدات والتمرين الكلي & 10\\
5 & أوزان الثلاثي المجرد & 12\\
6 & أمثال الثلاثي المجرد & 13\\
7 & أبواب الثلاثي المزيد والرباعي المجرد والمزيد & 14\\
8 & أوزان الثلاثي والرباعي والأمثال & 15\\
9 & الملحقات بالرباعي & 18\\
10 & افتعل تفعّل تفاعل & 19\\
11 & قال يقول (أجوف واوي) & 20\\
12 & باع يبيع (أجوف يائي) & 22\\
13 & أعان يعين (أجوف في أفعل) & 24\\
14 & دعا يدعو (ناقص واوي) & 26\\
15 & رمى يرمي (ناقص يائي) & 28\\
16 & آتى يؤتي (ناقص في أفعل) & 30\\
17 & مدّ يمدّ (مضاعف) & 32\\
18 & أحبّ يحبّ (مضاعف في أفعل) & 34\\
19 & اصطلاحات التعليل & 36\\
20 & الأفعال المعتلة وهي اثنتان ومائة & 37\\
\end{sarf}
\begin{sarf}{|c|>{\raggedleft\arraybackslash}p{.33\linewidth}|c|}
\thead{نمبر} & \thead{مضامين} & \thead{صفحه}\\\hline
21 & كل أفعال التعليل مع أسمائها وأوزانها & 55\\
22 & الأفعال المركبة & 56\\
23 & الناقص والأجوف في حالة الإعراب & 57\\
24 & ما يلحق مع الفعل واسمي الفاعل والمفعول & 58\\
25 & التصرف مع الضمائر اللازم والمتعدي & 59\\
26 & المصادر والصفة المشبهة واسم المبالغة & 60\\
27 & المذكر والمؤنث والمفرد والمجمع & 61\\
28 & كل الأفعال المستعملة & 61\\
29 & الاصطلاحات والتركيبات & 63\\
30 & العوامل ومتعلقاتها & 68\\
31 & تمرين كل الصرف & 74\\
32 & تمرين الصرف والتركيب والإعراب في الجمل & 75\\
33 & تمرين الصرف والتركيب والإعراب في الآيات & 76\\
34 & قوانين التعليل وغيره & 77\\
35 & طريق التعليم & 84\\
\end{sarf}
\end{tables}
